\section{Roadmap y fases}

\subsection{Descripcion}

Este documento define la secuencia de fases del trabajo documental y su relacion con el desarrollo del proyecto. Su objetivo es ordenar el avance del repositorio por reduccion real de ambiguedad y no por un simple orden superficial de carpetas.

\subsection{Alcance}

Aplica al repositorio documental y a su uso como base de trabajo para backend y frontend.

\subsection{Definiciones}

\begin{itemize}
  \item \textbf{Fase}: agrupacion de entregables con un objetivo comun.
  \item \textbf{Cierre funcional}: momento en que una seccion ya es suficientemente consistente para servir como referencia de trabajo.
\end{itemize}

\subsection{Reglas}

\begin{itemize}
  \item El trabajo documental se prioriza por reduccion de ambiguedad.
  \item Dominio y datos se cierran antes del detalle tecnico fino.
  \item Ninguna fase debe avanzar ignorando contradicciones abiertas de la fase anterior.
\end{itemize}

\subsubsection{Fases del roadmap documental}

\paragraph{Fase 1: nucleo del dominio}

\begin{itemize}
  \item objetivo: cerrar reglas de producto, estados, picks, locks, scoring, desempates y Edge Cases;
  \item entregables principales: \texttt{00\_Overview}, \texttt{01\_Producto\_y\_Reglas}, partes nucleares de \texttt{05\_Datos};
  \item resultado esperado: lenguaje comun y reglas del juego sin ambiguedad fuerte.
\end{itemize}

\paragraph{Fase 2: arquitectura y aterrizaje tecnico}

\begin{itemize}
  \item objetivo: traducir el dominio ya cerrado a arquitectura, modulos, fronteras y contratos conceptuales;
  \item entregables principales:
  \begin{itemize}
    \item \texttt{02\_Arquitectura};
    \item \texttt{03\_Backend};
    \item \texttt{04\_Frontend};
    \item \texttt{06\_Integraciones};
    \item \texttt{07\_Seguridad\_y\_Operacion}.
  \end{itemize}
  \item resultado esperado: base tecnica coherente para construir ambos repositorios de implementacion.
\end{itemize}

\paragraph{Fase 3: calidad, experiencia y gestion}

\begin{itemize}
  \item objetivo: cerrar testing, UI/UX, backlog documental, riesgos y ciclo de aceptacion;
  \item entregables principales: \texttt{08\_Testing}, \texttt{09\_UI\_UX}, \texttt{10\_Planeacion};
  \item resultado esperado: documentacion lista para trabajo iterativo, revision y salida controlada.
\end{itemize}

\subsubsection{Estado actual del roadmap}

Al momento de esta iteracion:

\begin{itemize}
  \item la Fase 1 esta sustancialmente cubierta;
  \item la Fase 2 esta mayormente cubierta en arquitectura, backend, frontend, integraciones y operacion;
  \item la Fase 3 esta en consolidacion documental y editorial.
\end{itemize}

\subsubsection{Siguiente tramo recomendado}

\begin{itemize}
  \item pasada de consistencia cruzada entre secciones;
  \item limpieza editorial y de maquetacion;
  \item compilacion del PDF maestro;
  \item refinamiento puntual de decisiones pendientes que aun bloquean frontend.
\end{itemize}

\subsection{Edge Cases}

\begin{itemize}
  \item Una fase puede requerir retrocesos puntuales si una decision abierta descubierta tarde afecta documentos ya avanzados.
  \item Un cierre parcial de fase es aceptable si los pendientes quedaron explicitamente registrados y no se presentan como decisiones ya resueltas.
\end{itemize}

\subsection{Invariantes}

\begin{itemize}
  \item No se adelanta detalle tecnico fino cuando el dominio aun esta abierto.
  \item Ninguna fase se considera cerrada si deja contradicciones normativas sin registrar.
  \item El roadmap documental debe poder explicar por que una seccion se trabajo antes que otra.
\end{itemize}

\subsection{Preguntas abiertas}

\begin{itemize}
  \item Confirmar si el proyecto quiere convertir estas fases en hitos formales por semana, sprint o solo mantenerlas como secuencia logica de trabajo.
\end{itemize}

\subsection{Decisiones pendientes}

\begin{itemize}
  \item \texttt{Decision pendiente} \texttt{No bloqueante}: decidir si el roadmap documental se seguira manteniendo como documento vivo o si solo reflejara grandes cortes de avance.
\end{itemize}
